% SPDX-FileCopyrightText: Copyright (c) 2016-2026 Objectionary.com
% SPDX-License-Identifier: MIT

\section{Well-formedness Rules}\label{sec:wellformedness}

\Cref{sec:syntax} described the per-line and multi-line shapes that EO admits.
This section gives the \emph{cross-line} rules that decide whether a
  syntactically-recognisable program is also legal:
  when a line may extend a previous one,
  when a name suffix is required,
  what an atom may contain,
  how comments attach,
  where blank lines belong,
  and how inline bindings may mix.
Together with the per-line syntax of \cref{sec:syntax},
  these rules form the complete well-formedness contract for EO source.

\subsection{Notation Consistency}\label{sec:wf:notation}

Each expression occupies a position in the \emph{openness lattice} ---
  one of \emph{open}, \emph{vertical-completed}, or \emph{horizontal-completed}.
The position is determined by the outer kind of the expression
  (\cref{app:outer-kinds}) and by how many additional lines have already
  attached.

A \emph{horizontally-completed} expression
  --- one whose outer kind lies in the set
\begin{equation*}
  \mathcal{H} = \{\,
    \mathtt{happlication},\,
    \mathtt{reversed\text{-}with\text{-}hargs},\,
    \mathtt{vmethod\text{-}with\text{-}hargs},\,
    \mathtt{inline\text{-}phi\text{-}formation}
    \,\} ---
\end{equation*}
admits neither deeper-indent children nor a same-indent \ff{.method}
  continuation.
A line that attempts to extend a horizontally-completed expression in either
  direction is rejected.

An \emph{open} expression --- one whose outer kind is
  \ff{head}, \ff{hmethod} (with 0 hargs), \ff{bare-formation},
  \ff{bare-reversed}, \ff{compact-tuple}, a \ff{vapplication} still receiving
  arguments, or a \ff{vmethod} chain whose latest \ff{.method} link has 0
  hargs --- may receive deeper-indent children but may not yet be wrapped by
  a same-indent \ff{.method}.

A \emph{vertical-completed} expression --- one whose body block has ended ---
  may no longer receive deeper-indent children, but a same-indent
  \ff{.method} continuation may still wrap it.
Wrapping transitions the expression into a \ff{vmethod} chain that is once
  again open for further continuation lines, until a \ff{.method}-with-hargs
  closes the chain horizontally.

\begin{ffcode}
foo                  // head, open
  a                  //   deeper child: foo becomes vapplication-in-progress
  b
.method > result     // same-indent .method after block closed:
                     // wraps as vmethod, naming-line carries the suffix

foo a b              // happlication, horizontally-completed
.method > result     // rejected: cannot wrap a horizontal application

io.stdout            // hmethod (0 hargs), open
  tt.sprintf         //   deeper child, fine

x.put 2              // happlication, horizontally-completed
  something          // rejected: closed for deeper-indent children
\end{ffcode}

A wrap of a vertical-completed multi-line expression into a fresh
  \ff{vmethod} chain is worked through in
  \cref{app:ex:vmethod-wrap}.

\subsection{Naming Inside Formations}\label{sec:wf:naming}

Every expression that is a \emph{plain child} of a formation must carry a
  name suffix on its \emph{naming line}.
A plain child is any child expression whose outer kind is not itself a
  formation-introducing form
  --- i.e., any \ff{head}, \ff{hmethod}, \ff{happlication},
  \ff{vapplication}, \ff{vmethod}, \ff{bare-reversed},
  \ff{reversed-with-hargs}, \ff{compact-tuple}, or \ff{text-block}.
A child whose outer kind is \ff{bare-formation} or
  \ff{inline-phi-formation}, or whose suffix is \ff{> name /sig} (an atom),
  is a \emph{master child}; the bracket marker makes its role self-evident,
  so a name suffix on a master child is optional.

The naming line for each outer kind --- the line on which the suffix must
  live --- is given in \cref{tab:naming-lines}.

\begin{table*}
\capt{Naming line per outer kind.}
\label{tab:naming-lines}
\small
\begin{tabularx}{\textwidth}{X X}
\toprule
Outer kind & Naming line \\
\midrule
Single-line expression (\ff{head}, \ff{hmethod}, \ff{happlication}, \ff{reversed-with-hargs}, \ff{inline-phi-formation})
  & The line itself. \\
\ff{vapplication}
  & The head line. \\
\ff{vmethod} (a chain of same-indent \ff{.method} continuations)
  & The \emph{last} \ff{.method} line of the chain. \\
\ff{vapplication} whose head is a \ff{vmethod}
  & The last \ff{.method} line (which is also the vapplication's head). \\
\ff{bare-formation} (with a body)
  & The \ff{[\ldots]} line. \\
\ff{bare-reversed} (vertical form)
  & The \ff{name.} line. \\
\ff{compact-tuple}
  & The head line carrying \ff{*N}. \\
\ff{text-block}
  & The closing-\ff{"""} line. \\
\bottomrule
\end{tabularx}
\end{table*}

Intermediate lines of a multi-line expression may carry their own optional
  names; these intermediate names bind sub-expressions and are independent
  of the outer expression's name.

\begin{ffcode}
[] > foo
  x.y.z > name        // single-line plain child, named on its line

  tmpdir              // intermediate, no name (optional)
  .tmpfile > file     // intermediate, named (optional)
  .deleted            // intermediate, no name
  .open > @           // naming line: outer vmethod is named @

  malloc.for > result // vapplication head, naming line
    0
    [m] >>
\end{ffcode}

Illegal:
\begin{ffcode}
[] > foo
  x.y.z               // rejected: plain child has no name

  tmpdir              // outer vmethod starts here
  .tmpfile
  .open               // rejected: vmethod's naming line carries no suffix
\end{ffcode}

\subsection{Atoms and Test Attributes}\label{sec:wf:atoms}

A formation whose suffix is \ff{> name /sig} is an \emph{atom}.
The signature \ff{sig} declares the atom's return type;
  the atom's behaviour is supplied by a host-language function via the
  \(\lambda\)-asset of the formation (\cref{sec:sem:atoms}).
A non-atom formation may contain any mix of plain children, master children,
  and test attributes; an atom may contain \emph{only} test attributes.
A plain child or master child appearing inside an atom is rejected.

A \emph{test attribute} is a child whose suffix is \ff{+> name}.
Test attributes are legal only at indent level 1 of a top-level object ---
  that is, as a direct child of a formation that itself sits at indent 0.
A \ff{+>} attribute at any other depth is rejected.
The \ff{name} of a test attribute must be a \ff{NAME} token;
  \ff{+> @} is rejected.

A nested atom --- an atom that does not sit at indent 0 --- inherits two
  restrictions.
Tests are admissible only as direct children of a top-level object, so a
  nested atom cannot hold tests;
  and an atom admits no plain or master children, so a nested atom cannot
  hold those either.
The body of a nested atom must therefore be \emph{empty}.

A nested atom inside another atom is rejected outright:
  the outer atom admits only test attributes, and an atom (a master child)
  is not a test attribute, so its presence violates the outer atom's body
  contract regardless of the inner atom's body shape.

\begin{ffcode}
[] > foo                // top-level formation
  bar > attr1           //   plain child
  [] > inner-formation  //   master child (regular)
  [] +> test-foo        //   test attribute, legal at indent 2

[] > foo /number        // top-level atom
  [] +> test-foo        //   only test attributes allowed inside

[] > foo
  [] > bar
    [] > baz /number    // nested atom at indent 4; legal but body must be empty
\end{ffcode}

Illegal:
\begin{ffcode}
[] > foo /number
  bar > attr            // rejected: atom cannot have regular children

[] > foo
  [] > inner
    [] +> test-inner    // rejected: tests must be at indent 2 of top-level
\end{ffcode}

A nested atom with an empty body is worked through in
  \cref{app:ex:nested-atom}.

\subsection{Comment Attachment}\label{sec:wf:comments}

A comment block (\cref{sec:syntax:comment}) attaches to the next named
  object at the same indent.
The attachment is structural: a comment block at indent \(K\) may attach
  only to a named object at indent \(K\), and there must be no blank line
  between the block and the object it documents.
A comment block followed by a blank line, by an object at a different
  indent, or by no named object at all, is \emph{dangling}
  and is rejected.

Two comment blocks at the same indent separated by one or more blank lines
  are not merged: the earlier block is dangling
  (no named object follows it directly),
  and only the later block attaches.
Inside an open formation body, a comment block is deferred until the next
  named child appears at that indent.

\begin{ffcode}
# Tests that while loop dataizes only the first cycle.
[] +> tests-while-dataizes    // comment attaches to this formation

# This comment leads nowhere.
                              // blank line breaks attachment
[] > foo                      // comment above is dangling -- rejected
\end{ffcode}

Further dangling-comment patterns (a comment block at a different
  indent from the following object, or two blocks separated by a
  blank line) are worked through in \cref{app:ex:dangling-comment}.

\subsection{Blank-line Policy}\label{sec:wf:blanks}

Blank lines have no indent and no kind of their own; their legality is
  decided by the lines they separate.

\begin{itemize}
  \item Inside a comment block: blank lines are forbidden.
  \item Between a comment block and the named object it documents:
    blank lines are forbidden.
  \item Before a master object (a formation, an atom, an inline-phi
    formation, or a \ff{+>} test attribute) ---
    or before its preceding comment block, if any ---
    zero or one blank line is permitted.
  \item Between two plain siblings: blank lines are forbidden.
  \item After the meta header: exactly one blank line separates metas from
    the rest of the program.
  \item At end-of-file: zero or one trailing blank line; more than one is
    an error.
\end{itemize}

\begin{ffcode}
+architect yegor256@gmail.com
+version 0.0.0
                              // exactly one blank after metas
# Top-level comment.
[] > foo
  a > attr1
  b > attr2                   // no blank between plain siblings
                              // blank legal before next master
  # Inner comment block.
  [] > inner
    body > x
\end{ffcode}

Illegal:
\begin{ffcode}
[] > foo
  a > attr1
                              // rejected: blank between plain siblings
  b > attr2
\end{ffcode}

\subsection{Bindings}\label{sec:wf:bindings}

An inline binding \ff{:label} or \ff{:N} (\cref{sec:syntax:binding}) names
  the attribute slot the argument attaches to.
Without a binding, arguments attach by positional index (\(\alpha_0\),
  \(\alpha_1\), \ldots).

Three constraints govern the use of bindings.

\paragraph{All-or-nothing.}
Within a same-indent group of application arguments, bindings must be
  either all present or all absent.
A mixed group is rejected.

\paragraph{Receiver exclusion.}
The receiver of a reversed-dispatch line is the dispatch target, not an
  argument, and cannot carry a binding.
The remaining arguments of the dispatch form an ordinary application-argument
  group, governed independently by the all-or-nothing rule.

\paragraph{Last-link-only on method chains.}
An inline binding on a method dispatch is permitted only when the method
  is the last link in its chain.
A binding on an intermediate chain link is rejected.

\begin{ffcode}
if. true x y               // legal: all unbound
if. true x:a y:b           // legal: receiver unbound, args all bound
foo a:1 b:2                // legal: all bound
foo a b                    // legal: all unbound
\end{ffcode}

Illegal:
\begin{ffcode}
foo a:1 b                  // rejected: mixed binding
if. true:a x y             // rejected: receiver carries binding
if. true x:a y             // rejected: args mixed
\end{ffcode}

The interaction between all-or-nothing binding and receiver
  exclusion in a reversed dispatch is worked through in
  \cref{app:ex:binding-chain}.
